\section{Introdução}
\label{sec:introducao}

A alocação de bens, contratos e recursos escassos quando o vendedor desconhece quanto cada comprador está disposto a pagar é um dos problemas mais importantes da economia e da engenharia. Na prática, leilões movimentam trilhões de dólares anualmente: governos leiloam frequências de telecomunicações (como o 5G), campos de petróleo e concessões de energia, enquanto plataformas digitais como o Google realizam bilhões de leilões instantâneos todos os dias para vender espaços de anúncios na internet \cite{klemperer2004,edelman2007,despotakis2021}.

Do ponto de vista da Teoria dos Jogos, um leilão difere dos jogos clássicos de informação completa (como o Dilema dos Prisioneiros ou o xadrez) porque envolve \textbf{informação incompleta}: cada participante sabe quanto o objeto vale para si, mas desconhece o quanto ele vale para os seus concorrentes. Até a década de 1960, não se sabia como calcular um Equilíbrio de Nash nesse cenário sem cair em uma regressão infinita de crenças. Esse problema foi resolvido pelos trabalhos clássicos de \cite{vickrey1961} e \cite{harsanyi1967} (laureados com o Prêmio Nobel), que modelaram os leilões como \textbf{jogos bayesianos} por meio dos conceitos de \textit{Tipo Privado} e de \textit{Prior Comum} \cite{gibbons1992}.

Neste artigo, apresentamos de forma construtiva e acessível a teoria dos leilões como jogos bayesianos. Na Seção~\ref{sec:conceito}, definimos sequencialmente os conceitos fundamentais (tipos, \textit{prior} comum, estratégias funcionais, Equilíbrio de Nash Bayesiano, valores privados vs. valores comuns, os quatro formatos clássicos de leilão, o Teorema da Equivalência de Receita e a Maldição do Vencedor). Na Seção~\ref{sec:modelo_resolvido}, resolvemos passo a passo o leilão de primeiro preço com dois participantes e valores uniformes em $[0,1]$, comparando sua receita com a do leilão de segundo preço. Na Seção~\ref{sec:caso_real}, analisamos e modelamos o caso real dos leilões de anúncios online (Google Ads e \textit{Ad Exchanges}). Por fim, a Seção~\ref{sec:ligacao_curso} conecta o tema à disciplina GED-62 e a Seção~\ref{sec:conclusao} apresenta as conclusões.
